% ============================================================
\chapter{Espaces Sobres et Locales}
\label{ch:sobres_locales}
% ============================================================

Et si l'on pouvait faire de la topologie sans jamais parler de points~? L'id\'ee semble absurde~: depuis Cantor et Hausdorff, un espace topologique est un ensemble de \emph{points} muni d'une collection d'ouverts. Pourtant, \`a partir des ann\'ees 1960, une tradition issue de la logique et de l'informatique th\'eorique --- port\'ee par des math\'ematiciens comme John Isbell et les th\'eoriciens des cat\'egories --- a montr\'e que les ouverts eux-m\^emes contiennent toute l'information. Le treillis des ouverts, vu comme structure alg\'ebrique, suffit \`a reconstruire la topologie. Les espaces o\`u cette reconstruction est fid\`ele sont les espaces \emph{sobres}~; la structure alg\'ebrique abstraite est une \emph{locale}.

Ce chapitre explore ce passage fascinant du concret (les points) \`a l'abstrait (les ouverts comme entit\'es autonomes), un voyage qui r\'ev\`ele des connexions profondes entre topologie, th\'eorie de l'ordre et logique.

\begin{intuition}
Les espaces sobres constituent la classe d'espaces $T_0$ la plus naturelle~: ce sont ceux
o\`u la topologie d\'etermine compl\`etement la structure ponctuelle, via la correspondance
entre points et ferm\'es irr\'eductibles. Les locales formalisent l'id\'ee d'une
<<~topologie sans points~>> et fournissent un cadre alg\'ebrique pour la topologie.
\end{intuition}

% ------------------------------------------------------------
\section{Ferm\'es irr\'eductibles}
% ------------------------------------------------------------

\begin{definition}[Ferm\'e irr\'eductible]
Un ferm\'e $F$ d'un espace topologique $X$ est \emph{irr\'eductible} si~:
\begin{enumerate}
  \item $F \neq \emptyset$;
  \item si $F = F_1 \cup F_2$ avec $F_1, F_2$ ferm\'es dans $F$, alors $F = F_1$ ou
    $F = F_2$.
\end{enumerate}
De mani\`ere \'equivalente, $F$ est irr\'eductible si et seulement si pour tous ouverts
$U_1, U_2$ de $X$ tels que $F \cap U_1 \neq \emptyset$ et $F \cap U_2 \neq \emptyset$,
on a $F \cap U_1 \cap U_2 \neq \emptyset$.
\end{definition}

\begin{proposition}[Propri\'et\'es des ferm\'es irr\'eductibles]
\begin{enumerate}
  \item L'adh\'erence d'un singleton $\overline{\{x\}}$ est toujours irr\'eductible.
  \item L'image continue d'un irr\'eductible est irr\'eductible.
  \item L'adh\'erence d'un irr\'eductible est irr\'eductible.
  \item Si $F$ est irr\'eductible et $F \cap G \neq \emptyset$ avec $G$ ouvert, alors
    $F \cap G$ est irr\'eductible dans $G$.
\end{enumerate}
\end{proposition}

\begin{proof}
(1) Si $\overline{\{x\}} = F_1 \cup F_2$ avec $F_i$ ferm\'es, alors $x \in F_1$ ou
$x \in F_2$, disons $x \in F_1$. Alors $\overline{\{x\}} \subseteq F_1$, donc
$F_1 = \overline{\{x\}}$.

(2) Soit $f\colon X \to Y$ continue et $F$ irr\'eductible dans $X$. Si
$f(F) \subseteq G_1 \cup G_2$ avec $G_i$ ferm\'es, alors $F \subseteq f^{-1}(G_1) \cup
f^{-1}(G_2)$ avec $f^{-1}(G_i)$ ferm\'es, donc $F \subseteq f^{-1}(G_i)$ pour un $i$,
d'o\`u $f(F) \subseteq G_i$. L'image $f(F)$ est irr\'eductible, et $\overline{f(F)}$
l'est aussi par~(3).

(3) Si $\overline{F} = G_1 \cup G_2$, alors $F \subseteq G_1 \cup G_2$, donc
$F = (F \cap G_1) \cup (F \cap G_2)$. Par irr\'eductibilit\'e, $F \subseteq G_i$ pour
un $i$, d'o\`u $\overline{F} \subseteq G_i$.
\end{proof}

% ------------------------------------------------------------
\section{Espaces sobres}
% ------------------------------------------------------------

\begin{definition}[Espace sobre]
Un espace topologique $X$ est \emph{sobre} si tout ferm\'e irr\'eductible poss\`ede un
unique point g\'en\'erique~: pour tout ferm\'e irr\'eductible $F$, il existe un unique
$x \in X$ tel que $\overline{\{x\}} = F$.
\end{definition}

\begin{theoreme}[Caract\'erisation des espaces sobres]
Un espace $X$ est sobre si et seulement si~:
\begin{enumerate}
  \item $X$ est $T_0$, et
  \item tout ferm\'e irr\'eductible a un point g\'en\'erique.
\end{enumerate}
\end{theoreme}

\begin{proof}
$(\Rightarrow)$ L'unicit\'e du point g\'en\'erique implique $T_0$~: si
$\overline{\{x\}} = \overline{\{y\}}$, alors $x$ et $y$ sont tous deux points
g\'en\'eriques du m\^eme ferm\'e irr\'eductible, donc $x = y$.

$(\Leftarrow)$ Il suffit de montrer l'unicit\'e. Si $\overline{\{x\}} = \overline{\{y\}} = F$,
alors comme $X$ est $T_0$, $x = y$.
\end{proof}

\begin{exemple}[Exemples d'espaces sobres]
\begin{enumerate}
  \item Tout espace de Hausdorff est sobre (les ferm\'es irr\'eductibles sont les
    singletons, car dans un espace $T_2$, tout ferm\'e irr\'eductible est un singleton).
  \item Le spectre $\mathrm{Spec}(A)$ d'un anneau commutatif est sobre (th\'eor\`eme
    classique de g\'eom\'etrie alg\'ebrique).
  \item L'espace de Sierpi\'nski $\mathbb{S}$ est sobre.
  \item Tout treillis complet avec la topologie de Scott est sobre.
\end{enumerate}
\end{exemple}

\begin{exemple}[Espace $T_0$ non sobre]
Soit $X = \N$ muni de la topologie dont les ouverts sont $\emptyset$ et les ensembles
cofinaux $\{n, n+1, n+2, \ldots\}$. L'ordre de sp\'ecialisation est l'ordre usuel sur
$\N$. Le ferm\'e $X$ lui-m\^eme est irr\'eductible (car tout ouvert non vide est cofinal
et deux ouverts non vides s'intersectent toujours). Mais $X = \overline{\{n\}}$ n'est
possible pour aucun $n$, car $\overline{\{n\}} = \{0, 1, \ldots, n\}$. Donc $X$ n'est
pas sobre.
\end{exemple}

% ------------------------------------------------------------
\section{La sobrification}
% ------------------------------------------------------------

\begin{definition}[Sobrification]
Soit $X$ un espace topologique. L'espace des \emph{ferm\'es irr\'eductibles} de $X$,
not\'e $X^s$ (ou $\mathrm{Sob}(X)$), est l'ensemble des ferm\'es irr\'eductibles de $X$,
muni de la topologie dont les ouverts sont~:
\[
  \tilde{U} = \{F \in X^s : F \cap U \neq \emptyset\}, \quad U \in \mathcal{O}(X).
\]
L'application $\eta_X\colon X \to X^s$ d\'efinie par $\eta_X(x) = \overline{\{x\}}$ est
continue.
\end{definition}

\begin{theoreme}[Propri\'et\'es de la sobrification]
\begin{enumerate}
  \item $X^s$ est sobre.
  \item L'application $U \mapsto \tilde{U}$ est un isomorphisme de cadres
    $\mathcal{O}(X) \cong \mathcal{O}(X^s)$.
  \item $\eta_X$ est universelle~: pour toute application continue $f\colon X \to Y$ avec
    $Y$ sobre, il existe une unique application continue $\tilde{f}\colon X^s \to Y$ telle
    que $f = \tilde{f} \circ \eta_X$.
  \item Le foncteur de sobrification $\mathrm{Sob}\colon \mathbf{Top} \to \mathbf{Sob}$
    est adjoint \`a gauche de l'inclusion $\mathbf{Sob} \hookrightarrow \mathbf{Top}$.
\end{enumerate}
\end{theoreme}

\begin{proof}
(1) Montrons que $X^s$ est $T_0$. Si $F_1 \neq F_2$ sont deux ferm\'es irr\'eductibles,
alors il existe un ouvert $U$ tel que, disons, $F_1 \cap U \neq \emptyset$ et
$F_2 \cap U = \emptyset$. Alors $F_1 \in \tilde{U}$ et $F_2 \notin \tilde{U}$.

Montrons que tout ferm\'e irr\'eductible de $X^s$ a un point g\'en\'erique. Soit
$\mathcal{F}$ un ferm\'e irr\'eductible de $X^s$. Posons
$F_0 = \overline{\bigcup_{F \in \mathcal{F}} F}$ dans $X$. On v\'erifie que $F_0$ est
un ferm\'e irr\'eductible de $X$, et que $\overline{\{F_0\}}^{X^s} = \mathcal{F}$.

(2) L'application $U \mapsto \tilde{U}$ pr\'eserve les r\'eunions quelconques et les
intersections finies par construction, et elle est bijective car $\tilde{U_1} = \tilde{U_2}$
implique, pour tout ferm\'e irr\'eductible $F$, que $F \cap U_1 \neq \emptyset \Leftrightarrow
F \cap U_2 \neq \emptyset$. Comme les singletons $\overline{\{x\}}$ sont des ferm\'es
irr\'eductibles, on obtient $U_1 = U_2$.

(3) Si $f\colon X \to Y$ est continue et $Y$ sobre, on d\'efinit
$\tilde{f}(F) = $ l'unique point g\'en\'erique de $\overline{f(F)}$ dans $Y$ (qui existe
car $f(F)$ est irr\'eductible par la proposition pr\'ec\'edente, et $Y$ est sobre). On
v\'erifie que $\tilde{f}$ est continue et unique.
\end{proof}

% ------------------------------------------------------------
\section{Cadres et locales}
% ------------------------------------------------------------

\begin{definition}[Cadre (\emph{Frame})]
Un \emph{cadre} est un treillis complet $L$ satisfaisant la loi de distributivit\'e infinie~:
\[
  a \wedge \bigvee_{i\in I} b_i = \bigvee_{i\in I}(a \wedge b_i)
\]
pour tout $a \in L$ et toute famille $(b_i)_{i\in I}$ dans $L$.
Un \emph{morphisme de cadres} $h\colon L \to M$ pr\'eserve les sups quelconques et les
infs finis (y compris $\top$ et $\bot$).
\end{definition}

\begin{definition}[Locale]
La cat\'egorie $\mathbf{Loc}$ des \emph{locales} est la cat\'egorie oppos\'ee de la
cat\'egorie des cadres~:
\[
  \mathbf{Loc} = \mathbf{Frm}^{\mathrm{op}}.
\]
Un objet de $\mathbf{Loc}$ est un cadre vu comme <<~espace sans points~>>, et un morphisme
de locales $f\colon L \to M$ est un morphisme de cadres $f^*\colon M \to L$.
\end{definition}

\begin{remarque}
Le renversement des fl\`eches est motiv\'e par l'analogie avec les espaces topologiques~:
une application continue $f\colon X \to Y$ induit $f^{-1}\colon \mathcal{O}(Y) \to
\mathcal{O}(X)$, qui va dans le sens oppos\'e.
\end{remarque}

% ------------------------------------------------------------
\section{Le foncteur $\mathrm{pt}$ et l'adjonction $\mathcal{O} \dashv \mathrm{pt}$}
% ------------------------------------------------------------

\begin{definition}[Points d'une locale]
Soit $L$ une locale (c'est-\`a-dire un cadre). Un \emph{point} de $L$ est un morphisme de
cadres $p\colon L \to \{0,1\}$ (o\`u $\{0,1\}$ est le cadre \`a deux \'el\'ements), ou de
fa\c{c}on \'equivalente, un \'el\'ement compl\`etement premier $q \in L$ ($q \neq \top$ et
$a \wedge b \leq q \Rightarrow a \leq q$ ou $b \leq q$).

L'espace des points $\mathrm{pt}(L)$ est l'ensemble des points de $L$ muni de la topologie
dont les ouverts sont les $\Sigma_a = \{p : p(a) = 1\}$ pour $a \in L$.
\end{definition}

\begin{theoreme}[Adjonction $\mathcal{O} \dashv \mathrm{pt}$]
Les foncteurs $\mathcal{O}\colon \mathbf{Top} \to \mathbf{Loc}$ et
$\mathrm{pt}\colon \mathbf{Loc} \to \mathbf{Top}$ forment une adjonction~:
\[
  \mathcal{O} \dashv \mathrm{pt}.
\]
L'unit\'e $\eta_X\colon X \to \mathrm{pt}(\mathcal{O}(X))$ envoie $x$ sur le point
$\eta_X(x)(U) = \begin{cases} 1 & \text{si } x \in U \\ 0 & \text{sinon}\end{cases}$.
\end{theoreme}

\begin{proof}
Pour un espace $X$ et une locale $L$, on doit \'etablir~:
\[
  \mathrm{Hom}_{\mathbf{Top}}(X, \mathrm{pt}(L)) \cong
  \mathrm{Hom}_{\mathbf{Loc}}(\mathcal{O}(X), L).
\]
Le membre droit est $\mathrm{Hom}_{\mathbf{Frm}}(L, \mathcal{O}(X))$. Une application
continue $f\colon X \to \mathrm{pt}(L)$ induit un morphisme de cadres
$L \to \mathcal{O}(X)$ envoyant $a \mapsto f^{-1}(\Sigma_a)$. R\'eciproquement, un
morphisme de cadres $h\colon L \to \mathcal{O}(X)$ induit une application continue
$X \to \mathrm{pt}(L)$ envoyant $x$ sur le point $a \mapsto [x \in h(a)]$.
\end{proof}

\begin{theoreme}[Caract\'erisation des espaces sobres via l'adjonction]
Un espace $X$ est sobre si et seulement si l'unit\'e $\eta_X\colon X \to
\mathrm{pt}(\mathcal{O}(X))$ est un hom\'eomorphisme.
\end{theoreme}

\begin{proof}
L'application $\eta_X$ est injective si et seulement si $X$ est $T_0$ (deux points
distincts donnent des ``points'' diff\'erents du cadre). Elle est surjective si et seulement
si tout point du cadre $\mathcal{O}(X)$ provient d'un point de $X$, ce qui \'equivaut
\`a l'existence d'un point g\'en\'erique pour chaque ferm\'e irr\'eductible. Enfin, $\eta_X$
est toujours un plongement topologique.
\end{proof}

% ------------------------------------------------------------
\section{Locales spatiales et cadres spatiaux}
% ------------------------------------------------------------

\begin{definition}[Locale spatiale]
Une locale $L$ est \emph{spatiale} si le morphisme canonique de cadres
$L \to \mathcal{O}(\mathrm{pt}(L))$ est un isomorphisme, c'est-\`a-dire si $L$ a
<<~assez de points~>>.
\end{definition}

\begin{theoreme}
Les restrictions de l'adjonction $\mathcal{O} \dashv \mathrm{pt}$ induisent une
\'equivalence de cat\'egories entre la cat\'egorie des espaces sobres et la cat\'egorie des
locales spatiales.
\end{theoreme}

\begin{proposition}[Exemples de locales non spatiales]
Il existe des cadres sans aucun point~: par exemple, le cadre des ouverts r\'eguliers
de $\R$ modulo un certain id\'eal est un cadre non trivial sans point.
\end{proposition}

% ------------------------------------------------------------
\section{Sous-locales et noyaux}
% ------------------------------------------------------------

\begin{definition}[Noyau]
Soit $L$ un cadre. Un \emph{noyau} (\emph{nucleus}) sur $L$ est une application
$j\colon L \to L$ telle que~:
\begin{enumerate}
  \item $a \leq j(a)$ (inflation);
  \item $j(j(a)) = j(a)$ (idempotence);
  \item $j(a \wedge b) = j(a) \wedge j(b)$ (pr\'eservation des infs finis).
\end{enumerate}
\end{definition}

\begin{theoreme}[Sous-locales via les noyaux]
Les sous-locales d'une locale $L$ correspondent bijectivement aux noyaux sur $L$. Si $j$
est un noyau, l'image $j(L) = \{a \in L : j(a) = a\}$ est un cadre (avec les m\^emes
infs que $L$ et les sups modifi\'es $\bigvee^j a_i = j(\bigvee a_i)$), et le morphisme
de cadres $L \to j(L)$, $a \mapsto j(a)$, d\'efinit un plongement de locales
$j(L) \hookrightarrow L$.
\end{theoreme}

\begin{proof}
V\'erifions que $j(L)$ est un cadre. Il est clairement ordonn\'e par l'ordre h\'erit\'e de
$L$. Pour les infs~: si $a, b \in j(L)$, alors $j(a \wedge b) = j(a) \wedge j(b) =
a \wedge b$, donc $a \wedge b \in j(L)$. Pour les sups~: $j(\bigvee a_i)$ est bien dans
$j(L)$ et c'est le plus petit \'el\'ement de $j(L)$ au-dessus de tous les $a_i$.
La distributivit\'e se v\'erifie en utilisant celle de $L$.
\end{proof}

% ------------------------------------------------------------
\section{Dualit\'e de Stone}
% ------------------------------------------------------------

\begin{definition}[Treillis distributif born\'e]
Un \emph{treillis distributif born\'e} est un treillis $D$ poss\'edant un plus petit
\'el\'ement $0$ et un plus grand \'el\'ement $1$, et satisfaisant
$a \wedge (b \vee c) = (a \wedge b) \vee (a \wedge c)$.
\end{definition}

\begin{definition}[Espace spectral]
Un espace topologique $X$ est \emph{spectral} s'il est~:
\begin{enumerate}
  \item $T_0$ et sobre;
  \item compact;
  \item l'ensemble des ouverts compacts est stable par intersection finie et forme une base
    de la topologie.
\end{enumerate}
\end{definition}

\begin{theoreme}[Dualit\'e de Stone pour les treillis distributifs]
Il y a une \'equivalence de cat\'egories entre la cat\'egorie des treillis distributifs
born\'es et la cat\'egorie des espaces spectraux (avec applications continues spectrales).
Le foncteur dans un sens associe \`a un treillis $D$ l'espace de ses filtres premiers, et
dans l'autre, \`a un espace spectral $X$ le treillis de ses ouverts compacts.
\end{theoreme}

% ------------------------------------------------------------
\section{Exercices}
% ------------------------------------------------------------

\begin{exercice}
Montrer que l'espace de Sierpi\'nski $\mathbb{S}$ est sobre et d\'ecrire ses ferm\'es
irr\'eductibles.
\end{exercice}

\begin{exercice}
Montrer qu'un espace est sobre si et seulement s'il est un $d$-espace et tout ferm\'e
irr\'eductible a un point g\'en\'erique.
\end{exercice}

\begin{exercice}
Soit $L$ un cadre. Montrer que $\mathrm{pt}(L)$ est toujours un espace sobre.
\end{exercice}

\begin{exercice}
V\'erifier que le cadre $\mathcal{O}(\R)$ est spatial (c'est-\`a-dire que $\R$ avec la
topologie usuelle donne une locale spatiale).
\end{exercice}

\begin{exercice}[Sous-locales ouvertes et ferm\'ees]
Soit $L$ un cadre et $a \in L$. Montrer que~:
\begin{enumerate}
  \item Le noyau $j_a(x) = a \vee x$ d\'efinit une sous-locale ferm\'ee.
  \item Le noyau $j^a(x) = a \Rightarrow x$ (o\`u $\Rightarrow$ est l'implication de
    Heyting) d\'efinit une sous-locale ouverte.
\end{enumerate}
\end{exercice}

\begin{exercice}
Montrer que la sobrification de $\N$ (avec la topologie des ouverts cofinaux) est
hom\'eomorphe \`a $\N \cup \{\infty\}$ avec la topologie d'Alexandrov de l'ordre usuel.
\end{exercice}

\begin{exercice}[Alg\`ebre de Heyting]
Montrer que tout cadre est une alg\`ebre de Heyting~: pour tous $a, b$, il existe un plus
grand \'el\'ement $c$ tel que $a \wedge c \leq b$, not\'e $a \Rightarrow b$. Montrer que
$a \Rightarrow b = \bigvee\{c : a \wedge c \leq b\}$.
\end{exercice}
